% Overleaf: set "Main document" to main.tex; compiler: pdfLaTeX (or LaTeX).
\documentclass[paper=a4, fontsize=11pt]{scrartcl}

\newcommand{\assignment}{CS 6120 Final Project Report --- GTC Networking Copilot (Retrieval-Augmented Matching)}
\newcommand{\duedate}{April 16, 2026}
% Shared preamble (loaded after \documentclass and \assignment / \duedate).
\usepackage[utf8]{inputenc}
\usepackage[T1]{fontenc}
\usepackage[margin=1in]{geometry}
\usepackage{booktabs}
\usepackage{array}
\usepackage{url}
\usepackage[hidelinks]{hyperref}

\hypersetup{
  pdftitle={\assignment},
  pdfauthor={Conference Matching Platform},
}

\title{\assignment}
\date{\duedate}


\author{%
  \textbf{Xian Cao, Xueming Tang, Chenyang Li} \\
  \href{https://github.com/xiancao2024/Conference-Matching-Platform}{GitHub repository} (add LDAP/e-mails per course instructions)
}

\begin{document}

\maketitle

\noindent\textbf{Repository:} \url{https://github.com/xiancao2024/Conference-Matching-Platform}

\section{Proposed Objective}

Build a \textbf{retrieval-augmented} system that helps NVIDIA GTC-style attendees answer \emph{who should I meet?} over a large roster: \textbf{retrieve} top attendee profiles with hybrid search, then \textbf{generate} short connection rationales using a \textbf{local LLM} conditioned only on retrieved fields, so the model does not rely on parametric knowledge of private attendees.

\section{Background and Motivation}

\subsection*{Heilmeier-style framing}

We follow the \href{https://www.darpa.mil/work-with-us/heilmeier-catechism}{Heilmeier Catechism} as a skeptical-audience check. \textbf{What are we trying to do?} Turn informal networking goals into ranked, explainable introductions. \textbf{How is it done today?} Spreadsheets, hallway luck, or generic chatbots without a grounded attendee index. \textbf{What is new?} Hybrid retrieval over 10k+ structured profiles plus \textbf{grounded} local-LLM explanations (not inventing people). \textbf{Who cares?} Attendees and organizers. \textbf{Risks?} Synthetic data is not real registration; CPU LLM latency; explanations must cite retrieved text only.

\subsection*{Impact}

Conference ROI is often capped by \textbf{search friction}. A deployable roster + retrieval + grounded generation path is a scalable alternative to fine-tuning on private data.

\subsection*{Related work (brief)}

Hybrid retrieval (sparse lexical signals + dense embeddings + rules) is standard in search and modern \textbf{RAG} systems \cite{lewis2020retrieval}. We adapt the same \emph{retrieve then generate} pattern: retrieve \textbf{entities} (attendees), then generate one-line reasons from \textbf{field bundles} passed to the LLM.

\subsection*{Data source, access, and statistics}

\textbf{Primary artifact:} wide-row CSV (\(\geq 10{,}000\) rows for scale), one row per attendee: name, e-mail, education, major, job title, experience, interests, registered GTC-style agenda titles, bio/snippet.

\textbf{How to obtain / reproduce:} generate locally at no API cost:
\begin{verbatim}
python3 scripts/generate_gtc_wide_csv.py --rows 10000 \
  -o data/gtc_generated_10k.csv
python3 -m conference_matching.gtc_import \
  --input data/gtc_generated_10k.csv -o data/conference_gtc.json
\end{verbatim}
\noindent The repository documents the schema; real organizer CSVs can use the same column aliases.

\textbf{Statistics:} \(N \approx 10^4\) attendees; each row carries multi-field text (interests + agenda pipe-list + bio). \textbf{Why retrieval is hard:} short noisy text, overlapping jargon (CUDA, LLMs), synonymy across agenda titles, and the need to rank \textbf{people} rather than duplicate session entities.

\begin{thebibliography}{9}
\bibitem{lewis2020retrieval}
P. Lewis et al.,
\textit{Retrieval-Augmented Generation for Knowledge-Intensive NLP Tasks}.
NeurIPS, 2020.
\end{thebibliography}

\section{Proposed Approach / Implementation Details}

\begin{itemize}
  \item \textbf{Pre-processing:} CSV \(\rightarrow\) normalized JSON (\texttt{conference\_matching/gtc\_import}); normalize headers via aliases; split multi-value interest/agenda fields; synthesize searchable bio text per attendee.
  \item \textbf{Retrieval algorithms:} hybrid scoring in \texttt{conference\_matching/engine.py}---lexical overlap, ontology concept expansion, structured boosts (role/intent/topic), optional dense similarity using \texttt{sentence-transformers/all-MiniLM-L6-v2} with FAISS \texttt{IndexFlatIP} on L2-normalized vectors (\texttt{data/embeddings.npy}, \texttt{data/faiss.index}).
  \item \textbf{Noise remedies:} multi-signal fusion (not embedding-only); light tokenization; people-only intent so session titles do not become separate match rows; optional stricter topical filters for people-queries.
  \item \textbf{Scale:} batch encoding (e.g.\ 256), cached embeddings, FAISS index reuse when entity count matches; Docker volume for \texttt{data/} on VMs.
  \item \textbf{Design decisions (``chunking'' analogy):} attendees are atomic units; agenda text stays embedded in each profile (no arbitrary chunk overlap); query is flattened to searchable text for embedding.
  \item \textbf{Validation / evaluation:} qualitative demo on 10k roster; latency checks on CPU VM; optional unit tests; compare hybrid vs.\ keyword behavior on representative queries; LLM outputs audited against retrieved fields (grounding).
  \item \textbf{LLM:} \textbf{Ollama} locally (e.g.\ \texttt{llama3.2:1b} in the bundled Docker image) for \texttt{llm\_reason} re-rank/explanation; retrieval and ranking run without the LLM.
\end{itemize}

\section{Submission Guidelines}

Submit the PDF to Gradescope as a \emph{single} team submission; list all members and LDAP/e-mails at the top of the document as required by the course. This write-up is kept near \textbf{two--three pages} when compiled. Container build/run: see \texttt{README.md} and \texttt{Dockerfile}.

\end{document}
